% ═══════════════════════════════════════════════════════════════════════════════
%  CAPÍTULO — Strategy
% ═══════════════════════════════════════════════════════════════════════════════

\chapter{Strategy}

% ─── Situación Problema ──────────────────────────────────────────────────────
\section{Situación Problema}

Airbnb aplica precios dinámicos a sus propiedades: el precio por noche
varía según el algoritmo de tarificación configurado por el anfitrión.
Algunos anfitriones optan por \emph{precios estacionales} —sube en
verano y festivos—, otros por \emph{precios basados en demanda} —ajuste
automático según ocupación en el área— y otros por \emph{descuentos por
estancia prolongada} —reducción progresiva para reservas de más de siete
noches—.

Si toda esta lógica reside en la clase \texttt{Propiedad} mediante una
cadena de condicionales, agregar una nueva estrategia de precios implica
modificar esa clase y el riesgo de que cambios en una estrategia afecten
accidentalmente a las otras. Además, los anfitriones pueden querer cambiar
su estrategia de precios en cualquier momento sin afectar la disponibilidad
de la propiedad.

% ─── Solución ────────────────────────────────────────────────────────────────
\section{Solución}

El patrón Strategy extrae cada algoritmo de tarificación a una clase
independiente que implementa la interfaz \texttt{EstrategiaPrecios} con
el método \texttt{calcularPrecio()}. La clase \texttt{Propiedad} mantiene
una referencia a la estrategia activa y la delega para calcular el precio
de cualquier noche consultada.

El anfitrión puede cambiar su estrategia en tiempo de ejecución sin más
que sustituir el objeto de estrategia. Cada algoritmo queda aislado,
es probable de forma independiente y puede incorporarse una nueva
estrategia sin modificar \texttt{Propiedad} ni las estrategias existentes.

% ─── Diagrama de clases ─────────────────────────────────────────────────────
\begin{figure}[H]
    \centering
    % \includegraphics[width=\textwidth]{imagenes/abstract_factory_diagrama.png}
    \caption{Diagrama de clases del patrón Abstract Factory}
\end{figure}


% ─── Roles de las Clases ────────────────────────────────────────────────────
\section{Roles de las Clases}

% Explicar la responsabilidad de cada clase/interfaz

\begin{itemize}
    \item \textbf{AbstractFactory:}
    \item \textbf{ConcreteFactory:}
    \item \textbf{AbstractProduct:}
    \item \textbf{ConcreteProduct:}
    \item \textbf{Client:}
\end{itemize}


% ─── Main ───────────────────────────────────────────────────────────────────
\section{Main}

% Explicar qué realiza el programa principal

\begin{lstlisting}[language=Java, caption=NombreClase]
 
\end{lstlisting}


% ─── Salida del Main ────────────────────────────────────────────────────────
\section{Salida del Main}

\begin{lstlisting}[language=Java, caption=NombreClase]
 
\end{lstlisting}


% ─── Clases ─────────────────────────────────────────────────────────────────
\section{Clases}

% ─── Clase 1 ────────────────────────────────────────────────────────────────
\subsection{NombreClase}

\begin{lstlisting}[language=Java, caption=NombreClase]
 
\end{lstlisting}


% ─── Clase 2 ────────────────────────────────────────────────────────────────
\subsection{NombreClase}

\begin{lstlisting}[language=Java, caption=NombreClase]
 
\end{lstlisting}


% ─── Clase 3 ────────────────────────────────────────────────────────────────
\subsection{NombreClase}

\begin{lstlisting}[language=Java, caption=NombreClase]
 
\end{lstlisting}


% ─── Conclusiones ───────────────────────────────────────────────────────────
\section{Conclusiones}

% Explicar:
% - Ventajas del patrón
% - Cuándo usarlo
% - Beneficios obtenidos
% - Posibles desventajas
% - Relación con principios SOLID